\documentclass[11pt]{article}

\usepackage[margin=1in]{geometry}
\usepackage{amsmath}
\usepackage{array}
\usepackage{booktabs}
\usepackage{longtable}
\usepackage[T1]{fontenc}
\usepackage{lmodern}
\usepackage[hidelinks]{hyperref}

\title{Power Flow Solver Study Plan}
\author{Load Flow Lite}
\date{\today}

\newcommand{\code}[1]{\texttt{#1}}
\newcommand{\pu}{\mathrm{p.u.}}

\emergencystretch=2em
\sloppy

\begin{document}

\maketitle
\tableofcontents
\newpage

\section{Purpose}

This is a guided plan for understanding how this repository actually solves the
AC power flow problem. Treat it as a lab, not a checklist: read a small amount
of code, predict what it should do, run a focused test, then explain the result
back in your own words.

The shortest path through the solver is:

\begin{enumerate}
  \item Build the network admittance matrix, \code{Ybus}.
  \item Convert a voltage guess into calculated bus injections.
  \item Compare calculated injections with specified injections.
  \item Assemble the Newton--Raphson Jacobian.
  \item Solve a linear correction system.
  \item Update voltage angles and magnitudes.
  \item Repeat until mismatch is below tolerance.
  \item Compute branch flows from the solved voltages.
\end{enumerate}

\section{Learning Outcomes}

By the end, you should be able to:

\begin{itemize}
  \item Explain which variables are solved for at slack, PV, and PQ buses.
  \item Trace how a \code{Case} becomes a \code{Ybus} matrix.
  \item Derive $S = V \cdot \operatorname{conj}(Y_{\mathrm{bus}} V)$ from
  voltage and current.
  \item Build the mismatch vector in the same order as the code.
  \item Map every block of the Jacobian to its rows, columns, and bus types.
  \item Hand-compute one Newton step for a two-bus slack/PQ system.
  \item Explain how solved bus voltages become branch powers and losses.
  \item Identify which files are core numerical code and which are wrappers.
\end{itemize}

\section{How to Use This Plan}

Make two passes.

Pass 1 is top-down. Start with \code{solve\_power\_flow} so you know the whole
control flow before studying formulas.

Pass 2 is bottom-up. Study \code{build\_ybus}, \code{power\_injections},
\code{\_mismatch}, and \code{\_jacobian} in detail, then return to
\code{solve\_power\_flow}.

For each session:

\begin{enumerate}
  \item Read the listed functions.
  \item Answer the checkpoint questions without looking ahead.
  \item Run the focused test command.
  \item Write down one sentence explaining what the code is doing
  mathematically.
\end{enumerate}

The commands below assume the project virtual environment from the README:

\begin{verbatim}
python3 -m venv .venv
./.venv/bin/pip install numpy pytest
\end{verbatim}

\section{Function Map}

Use this as the map from concepts to code.

\renewcommand{\arraystretch}{1.18}
\begin{longtable}{p{0.19\textwidth}p{0.30\textwidth}p{0.41\textwidth}}
\toprule
Concept & Main code & Why it matters \\
\midrule
\endhead
Case validation &
\href{../src/powerflow/case.py\#L63}{\code{validate\_case}} &
Rejects invalid topology before Newton starts. \\
Bus indexing &
\href{../src/powerflow/case.py\#L58}{\code{bus\_index}} &
Converts external bus IDs to array indices. \\
Network model &
\href{../src/powerflow/ybus.py\#L10}{\code{build\_ybus}} &
Builds the admittance matrix used in every injection calculation. \\
Complex voltages &
\href{../src/powerflow/flows.py\#L35}{\code{complex\_voltages}} &
Converts magnitudes and angles into complex phasors. \\
Calculated injections &
\href{../src/powerflow/solver.py\#L36}{\code{power\_injections}} &
Computes \code{P\_calc} and \code{Q\_calc} from the current voltage guess. \\
Specified injections &
\href{../src/powerflow/solver.py\#L30}{\code{specified\_injections}} &
Applies the repo sign convention: generation minus load. \\
Bus type handling &
\href{../src/powerflow/solver.py\#L46}{\code{\_bus\_type\_indices}} &
Determines which equations and unknowns exist. \\
Initial guess &
\href{../src/powerflow/solver.py\#L53}{\code{\_initial\_voltage}} &
Creates the first Newton state. \\
Mismatch vector &
\href{../src/powerflow/solver.py\#L62}{\code{\_mismatch}} &
Defines the residual that Newton drives to zero. \\
Jacobian &
\href{../src/powerflow/solver.py\#L77}{\code{\_jacobian}} &
Linearizes the AC power-flow equations. \\
Newton loop &
\href{../src/powerflow/solver.py\#L139}{\code{solve\_power\_flow}} &
Orchestrates the full nonlinear solve. \\
Result assembly &
\href{../src/powerflow/solver.py\#L195}{\code{\_build\_result}} &
Converts the converged state into useful outputs. \\
Branch flows &
\href{../src/powerflow/flows.py\#L39}{\code{compute\_branch\_flows}} &
Computes line powers and losses after convergence. \\
\bottomrule
\end{longtable}

\section{Glossary}

\begin{description}
  \item[\code{Ybus}.] The bus admittance matrix. It maps complex bus voltages
  to complex injected currents: \code{I = Ybus @ V}.
  \item[\code{V}.] Complex bus voltage phasor. In code it is built from
  \code{v\_magnitude} and \code{v\_angle}.
  \item[\code{theta}.] Voltage angle in radians.
  \item[\code{P}, \code{Q}.] Real and reactive power injection. Positive means
  net injection into the network; negative means net load.
  \item[\code{S}.] Complex power, $S = P + jQ$.
  \item[\code{G}, \code{B}.] Real and imaginary parts of \code{Ybus}, so
  $Y_{\mathrm{bus}} = G + jB$.
  \item[Slack bus.] Fixed voltage magnitude and fixed angle. Its real and
  reactive generation are solved outputs.
  \item[PV bus.] Fixed real-power injection and voltage magnitude. Its angle is
  solved, and its reactive generation is a solved output.
  \item[PQ bus.] Fixed real and reactive injection. Its voltage angle and
  magnitude are solved.
  \item[Mismatch.] Specified injection minus calculated injection.
  \item[Jacobian.] The matrix of partial derivatives used to turn a nonlinear
  mismatch into a linear correction step.
  \item[Per unit.] Normalized power-system units. The repo stores powers in
  p.u.; conversion to MW/MVAr uses \code{base\_mva}.
\end{description}

\section{Session 1: Get the Whole Solver in Your Head}

Read:

\begin{itemize}
  \item \href{../src/powerflow/solver.py\#L139}{\code{solve\_power\_flow}}
  \item \href{../src/powerflow/solver.py\#L195}{\code{\_build\_result}}
\end{itemize}

Focus on the loop shape:

\begin{verbatim}
build Ybus
choose unknowns
initialize voltage state
repeat:
    compute mismatch
    stop if mismatch is small
    build Jacobian
    solve J * step = mismatch
    update angles and magnitudes
\end{verbatim}

Checkpoint questions:

\begin{itemize}
  \item Where is \code{Ybus} built?
  \item Where does the code stop on convergence?
  \item What arrays are updated during each Newton step?
  \item Why does non-convergence raise instead of returning a partial result?
\end{itemize}

Run:

\begin{verbatim}
./.venv/bin/python -m pytest \
  tests/test_solver.py::test_two_bus_power_flow_converges_and_balances_power -q
\end{verbatim}

Done criteria:

\begin{itemize}
  \item You can describe \code{solve\_power\_flow} without mentioning any
  wrapper code.
  \item You can point to the exact lines that call \code{\_mismatch},
  \code{\_jacobian}, and \code{np.linalg.solve}.
\end{itemize}

\section{Session 2: Understand Bus Types and the State Vector}

Read:

\begin{itemize}
  \item \href{../src/powerflow/solver.py\#L46}{\code{\_bus\_type\_indices}}
  \item \href{../src/powerflow/solver.py\#L53}{\code{\_initial\_voltage}}
  \item \href{../src/powerflow/case.py\#L63}{\code{validate\_case}}
\end{itemize}

The Newton state is not all bus voltages. It is:

\begin{itemize}
  \item Voltage angles for PV and PQ buses.
  \item Voltage magnitudes for PQ buses only.
\end{itemize}

The slack bus angle and magnitude stay fixed. PV bus voltage magnitudes also
stay fixed, while their reactive generation is solved after convergence.

Checkpoint questions:

\begin{itemize}
  \item What would go wrong if the slack angle were included as an unknown?
  \item Why does a PV bus solve for angle but not voltage magnitude?
  \item Why does a PQ bus solve for both angle and magnitude?
  \item For a system with 1 slack, 2 PV buses, and 6 PQ buses, how many Newton
  unknowns are there?
\end{itemize}

Run:

\begin{verbatim}
./.venv/bin/python -m pytest \
  tests/test_solver.py::test_pv_bus_keeps_voltage_magnitude_fixed -q
\end{verbatim}

Done criteria:

\begin{itemize}
  \item You can construct the state vector for any mix of slack, PV, and PQ
  buses.
  \item You can explain why PV-bus \code{q\_gen} is not a fixed input to the
  mismatch.
\end{itemize}

\section{Session 3: Build \code{Ybus}}

Read:

\begin{itemize}
  \item \href{../src/powerflow/ybus.py\#L10}{\code{build\_ybus}}
  \item \href{../src/powerflow/flows.py\#L23}{\code{\_branch\_admittance\_terms}}
\end{itemize}

\code{Ybus} is the electrical network model. Each branch contributes
self-admittance terms to the diagonal and mutual admittance terms to the
off-diagonal entries. Line charging, transformer tap ratios, phase shifts, and
bus shunts all enter here.

Checkpoint questions:

\begin{itemize}
  \item Why does a simple series line add \code{+y} to both diagonals and
  \code{-y} to both off-diagonals?
  \item Why is line charging split between the two ends?
  \item What changes when \code{tap\_ratio != 1.0}?
  \item Why can a phase shifter make \code{Ybus[from, to]} differ from
  \code{Ybus[to, from]}?
\end{itemize}

Run:

\begin{verbatim}
./.venv/bin/python -m pytest tests/test_ybus.py -q
\end{verbatim}

Done criteria:

\begin{itemize}
  \item You can manually stamp a two-bus branch into a 2x2 \code{Ybus}.
  \item You can explain why \code{build\_ybus} and
  \code{\_branch\_admittance\_terms} must use the same branch model.
\end{itemize}

\section{Session 4: Convert Voltages Into Power Injections}

Read:

\begin{itemize}
  \item \href{../src/powerflow/flows.py\#L35}{\code{complex\_voltages}}
  \item \href{../src/powerflow/solver.py\#L36}{\code{power\_injections}}
\end{itemize}

The key formula is:

\begin{align*}
I &= Y_{\mathrm{bus}} V, \\
S &= V \cdot \operatorname{conj}(I), \\
P_{\mathrm{calc}} &= \operatorname{real}(S), \\
Q_{\mathrm{calc}} &= \operatorname{imag}(S).
\end{align*}

This is where the current voltage guess becomes calculated real and reactive
power injection at each bus.

Checkpoint questions:

\begin{itemize}
  \item What are the shapes of \code{V}, \code{Ybus @ V}, and \code{S}?
  \item Why is the current conjugated in the complex-power formula?
  \item What does a negative \code{P\_calc} mean?
  \item What does a negative \code{Q\_calc} mean?
\end{itemize}

Run:

\begin{verbatim}
./.venv/bin/python -m pytest \
  tests/test_solver.py -k power_injection -q
\end{verbatim}

Done criteria:

\begin{itemize}
  \item You can compute \code{S = V * conj(Ybus @ V)} for a two-bus example.
  \item You can explain why this function knows nothing about slack, PV, or PQ
  bus types.
\end{itemize}

\section{Session 5: Build the Mismatch Vector}

Read:

\begin{itemize}
  \item \href{../src/powerflow/solver.py\#L30}{\code{specified\_injections}}
  \item \href{../src/powerflow/solver.py\#L62}{\code{\_mismatch}}
\end{itemize}

The sign convention is:

\begin{verbatim}
P_spec = p_gen - p_load
Q_spec = q_gen - q_load
mismatch = specified - calculated
\end{verbatim}

The mismatch vector is ordered as:

\begin{enumerate}
  \item Real-power mismatches for PV and PQ buses.
  \item Reactive-power mismatches for PQ buses.
\end{enumerate}

Checkpoint questions:

\begin{itemize}
  \item Why is there no real-power mismatch equation for the slack bus?
  \item Why is there no reactive-power mismatch equation for PV buses?
  \item How does the mismatch ordering match the Newton update vector?
  \item What would change if mismatch were defined as calculated minus
  specified?
\end{itemize}

Run:

\begin{verbatim}
./.venv/bin/python -m pytest \
  tests/test_solver.py::test_two_bus_power_flow_converges_and_balances_power -q
\end{verbatim}

Done criteria:

\begin{itemize}
  \item You can write the mismatch vector shape for a given set of bus types.
  \item You can explain the sign of the update in \code{solve\_power\_flow}.
\end{itemize}

\section{Session 6: Spend the Most Time on the Jacobian}

Read:

\begin{itemize}
  \item \href{../src/powerflow/solver.py\#L77}{\code{\_jacobian}}
\end{itemize}

This is the densest numerical function in the repo. It assembles the analytic
polar AC power-flow Jacobian:

\[
\begin{bmatrix}
\partial P_{\mathrm{calc}} / \partial \theta &
\partial P_{\mathrm{calc}} / \partial V \\
\partial Q_{\mathrm{calc}} / \partial \theta &
\partial Q_{\mathrm{calc}} / \partial V
\end{bmatrix}
\]

Rows match equations:

\begin{itemize}
  \item \code{P} rows for PV and PQ buses.
  \item \code{Q} rows for PQ buses.
\end{itemize}

Columns match unknowns:

\begin{itemize}
  \item Angle columns for PV and PQ buses.
  \item Voltage-magnitude columns for PQ buses.
\end{itemize}

The implementation uses derivatives of calculated injections. Since mismatch is
specified minus calculated, the Newton correction is applied consistently by
solving:

\begin{align*}
J \cdot \mathrm{step} &= \mathrm{mismatch}, \\
x &\leftarrow x + \mathrm{step}.
\end{align*}

Checkpoint questions:

\begin{itemize}
  \item Which nested loop fills \code{dP/dtheta}?
  \item Which nested loop fills \code{dP/dV}?
  \item Which nested loop fills \code{dQ/dtheta}?
  \item Which nested loop fills \code{dQ/dV}?
  \item Why do diagonal and off-diagonal formulas differ?
  \item How do \code{g = ybus.real} and \code{b = ybus.imag} enter the
  formulas?
\end{itemize}

Run:

\begin{verbatim}
./.venv/bin/python -m pytest tests/test_solver.py -q
\end{verbatim}

Done criteria:

\begin{itemize}
  \item You can label each quadrant of the Jacobian in the code.
  \item For a 1 slack, 1 PV, 2 PQ system, you can state the Jacobian dimensions
  and row/column ordering.
\end{itemize}

\section{Session 7: Work One Newton Step by Hand}

Use the two-bus case from
\href{../tests/test_solver.py\#L38}{\code{test\_two\_bus\_power\_flow\_converges\_and\_balances\_power}}:

\begin{verbatim}
bus 1: slack, V = 1.0, theta = 0
bus 2: PQ load, P_load = 0.4, Q_load = 0.2
branch: r = 0.02, x = 0.04, b = 0.0
flat start: V1 = 1.0 angle 0, V2 = 1.0 angle 0
\end{verbatim}

First compute the series admittance:

\[
y = \frac{1}{0.02 + j0.04} = 10 - j20.
\]

So the flat-start \code{Ybus} is:

\[
Y_{\mathrm{bus}} =
\begin{bmatrix}
10 - j20 & -10 + j20 \\
-10 + j20 & 10 - j20
\end{bmatrix}.
\]

At flat voltage, \code{Ybus @ V = 0}, so:

\begin{align*}
P_{\mathrm{calc},2} &= 0, &
Q_{\mathrm{calc},2} &= 0, \\
P_{\mathrm{spec},2} &= 0 - 0.4 = -0.4, &
Q_{\mathrm{spec},2} &= 0 - 0.2 = -0.2.
\end{align*}

Therefore:

\[
\mathrm{mismatch} =
\begin{bmatrix}
-0.4 \\
-0.2
\end{bmatrix}.
\]

For bus 2 at the flat start:

\[
G_{22} = 10,\quad B_{22} = -20,\quad V_2 = 1,\quad
P_{\mathrm{calc},2} = 0,\quad Q_{\mathrm{calc},2} = 0.
\]

The 2x2 Jacobian is:

\begin{align*}
\frac{\partial P}{\partial \theta}
  &= -Q_{\mathrm{calc},2} - B_{22} V_2^2 = 20, \\
\frac{\partial P}{\partial V}
  &= P_{\mathrm{calc},2} / V_2 + G_{22} V_2 = 10, \\
\frac{\partial Q}{\partial \theta}
  &= P_{\mathrm{calc},2} - G_{22} V_2^2 = -10, \\
\frac{\partial Q}{\partial V}
  &= Q_{\mathrm{calc},2} / V_2 - B_{22} V_2 = 20.
\end{align*}

\[
J =
\begin{bmatrix}
20 & 10 \\
-10 & 20
\end{bmatrix}.
\]

Solve:

\[
J \cdot \mathrm{step} =
\begin{bmatrix}
-0.4 \\
-0.2
\end{bmatrix},
\qquad
\mathrm{step} =
\begin{bmatrix}
-0.012 \\
-0.016
\end{bmatrix}.
\]

So the first update is approximately:

\begin{align*}
\theta_2 &= 0.0 - 0.012 = -0.012\ \mathrm{radians}, \\
V_2 &= 1.0 - 0.016 = 0.984\ \pu.
\end{align*}

Checkpoint questions:

\begin{itemize}
  \item Why is the first angle update negative?
  \item Why does the first voltage-magnitude update reduce \code{V2}?
  \item Which two entries of the state vector exist in this two-bus case?
  \item Which entries would disappear if bus 2 were PV instead of PQ?
\end{itemize}

Done criteria:

\begin{itemize}
  \item You can reproduce the first Newton step without reading the code.
  \item You can point from each hand-computed quantity to the line that computes
  it.
\end{itemize}

\section{Session 8: Interpret the Solved Result}

Read:

\begin{itemize}
  \item \href{../src/powerflow/solver.py\#L195}{\code{\_build\_result}}
  \item \href{../src/powerflow/flows.py\#L39}{\code{compute\_branch\_flows}}
  \item \href{../src/powerflow/solver.py\#L221}{\code{bus\_power\_balance\_residuals}}
\end{itemize}

After convergence, the code computes solved generation and branch flows. This is
post-solution interpretation, not part of Newton convergence.

Checkpoint questions:

\begin{itemize}
  \item Why is slack-bus generation computed after the solve?
  \item Why is PV-bus reactive generation computed after the solve?
  \item How are branch from-end and to-end currents computed?
  \item Why are branch losses the sum of the two terminal complex powers?
\end{itemize}

Run:

\begin{verbatim}
./.venv/bin/python -m pytest tests/test_flows.py -q
\end{verbatim}

Done criteria:

\begin{itemize}
  \item You can explain how solved bus voltages become branch \code{P/Q} flows.
  \item You can distinguish mismatch convergence from post-solve reporting.
\end{itemize}

\section{Session 9: Validate the Complete Implementation}

Run the focused tests first:

\begin{verbatim}
./.venv/bin/python -m pytest \
  tests/test_ybus.py tests/test_solver.py tests/test_flows.py -q
\end{verbatim}

Then run the broader suite:

\begin{verbatim}
./.venv/bin/python -m pytest -q
\end{verbatim}

The small tests are best for understanding individual equations. The IEEE tests
are best for confirming that the complete implementation matches static
MATPOWER fixtures when those fixtures are present.

Done criteria:

\begin{itemize}
  \item You know which test protects each concept.
  \item You can explain why a solver can pass a two-bus test but still need IEEE
  case validation.
\end{itemize}

\section{Common Confusions}

\begin{itemize}
  \item Slack-bus \code{p\_gen} and \code{q\_gen} in the input are not fixed
  mismatch equations; the solved slack injection is computed from the final
  network state.
  \item PV-bus \code{q\_gen} is solved, not fixed. This repo does not implement
  reactive power limits or PV-to-PQ conversion.
  \item PQ-bus voltage magnitude is unknown, even if the input provides an
  initial value.
  \item The mismatch excludes slack equations because the slack bus absorbs the
  active-power residual and fixes the angle reference.
  \item The mismatch excludes PV reactive equations because PV voltage magnitude
  is fixed and reactive output is allowed to move.
  \item \code{power\_injections} computes injections for every bus, but
  \code{\_mismatch} selects only the equations that belong in the Newton solve.
  \item \code{build\_ybus} builds a global network model;
  \code{compute\_branch\_flows} uses the same branch model locally for each line
  after convergence.
  \item The code stores angles in radians. Some fixtures and reports may display
  degrees.
  \item The code solves \code{J * step = mismatch} because the implementation
  uses the Jacobian of calculated injections together with a
  specified-minus-calculated mismatch convention.
\end{itemize}

\section{Optional Theory Backup}

If you want more context while reading the code, use:

\begin{itemize}
  \item \href{ieee9_power_flow_lecture_note.tex}{\code{docs/ieee9\_power\_flow\_lecture\_note.tex}},
  especially the Newton--Raphson appendix.
  \item \href{powerflow_v0_prompt.md}{\code{docs/powerflow\_v0\_prompt.md}},
  which describes the original implementation goals and constraints.
\end{itemize}

Use these as references only after you have tried to trace the code path. The
fastest learning loop is still: code, tiny example, test, explanation.

\section{Files to Delay Until Later}

These files are useful, but they are not the best starting point for
understanding the numerical solve:

\begin{itemize}
  \item \href{../src/powerflow/io.py}{\code{src/powerflow/io.py}}: case loading.
  \item \href{../src/powerflow/case.py}{\code{src/powerflow/case.py}}: data
  structures and validation.
  \item \href{../src/powerflow/contingency.py}{\code{src/powerflow/contingency.py}}:
  repeated solves for outage analysis.
  \item \href{../scripts/run_case9_n1.py}{\code{scripts/run\_case9\_n1.py}}:
  command-line report wrapper.
\end{itemize}

Read these after the Newton--Raphson path is clear.

\section{Final Self-Test}

You are done with the core study plan when you can answer these without looking
at the code:

\begin{itemize}
  \item What is the Newton state vector for a network with 1 slack, 2 PV, and 6
  PQ buses?
  \item Why is the slack bus excluded from both \code{P} and \code{Q} mismatch
  equations?
  \item Why is a PV bus included in \code{P} mismatch but excluded from
  \code{Q} mismatch?
  \item What does \code{S = V * conj(Ybus @ V)} compute?
  \item How is \code{P\_spec} different from \code{P\_calc}?
  \item What are the four Jacobian blocks?
  \item Why does the code update only \code{v\_angle[angle\_buses]} and
  \code{v\_magnitude[pq\_buses]}?
  \item How do branch flows differ from bus injections?
  \item Which function would you inspect first if voltages fail to converge?
  \item Which function would you inspect first if voltages converge but branch
  losses look wrong?
\end{itemize}

\end{document}
